\begin{center}\textls[100]{d. Ternäre quadratische Formen.\\Der VI. Abschnitt der D. A.}\end{center}
\begin{center}\bfseries 11.\end{center}
\noindent\hspace*{\parindent}Noch tiefer dringend aber ist das von \textsc{Gauss} erkannte Mittel, um die durch den zuvor angeführten Satz nur beschränkte Anzahl \(\gamma\) der Geschlechter genau zu bestimmen. Es ist die Theorie der ternären quadratischen Formen
\[
f=ax^2+a'x'^2+a''x''^2+2bx'x''+2b'x''x+2b''xx',
\]
für welche der Ausdruck
\[
\Delta=ab^2+a'b'^2+a''b''^2-aa'a''-2bb'b''
\]
die Determinante der Form genannt wird. An sich zwar lag der verallgemeinernde Fortschritt von den binären quadratischen Formen zu den quadratischen Formen mit mehreren Unbestimmten oder auch zu Formen höheren Grades nahe genug. Immerhin hat \textsc{Gauss} mit dem am 14. Februar 1799 (Nr. 96 T.) begonnenen erfolgreichen Beschreiten dieser Richtung der weiteren Forschung ein schier unermessliches Feld gewiesen, auf welchem dann später die Bemühungen von \textsc{Eisenstein}\footnote{\textsc{Crelles} Journal für Mathematik \textbf{35} (1847), S. 117, Gesammelte Abhandlungen, S. 177; Berliner Monatsberichte 1852, S. 356.}, von H. St. \textsc{Smith}\begingroup\renewcommand{\thefootnote}{**)}\footnote{London Transactions \textbf{157} (1867), S. 255, Papers I, S. 455. Mém. Sav. Étrang. Paris \textsuperscript{(2)} \textbf{29} (1887) Nr. 1, S. 55, Papers II, S. 677.}\endgroup\ und \textsc{Minkowski}\begingroup\renewcommand{\thefootnote}{***)}\footnote{Mém. Sav. Étrang. Paris \textsuperscript{(2)} \textbf{29} (1887) Nr. 2, S. 159, 164.}\endgroup\ u. A. den reichsten Ertrag geliefert haben. Für seinen Zweck bedurfte \textsc{Gauss} nur der ersten Elemente der Theorie. Diese laufen, obwohl sie mannigfaltigere
\newpage\SourceLeaf{33}\ArithmeticHead{26}{L}{BACHMANN, ÜBER GAUSS’ ZAHLENTHEORETISCHE ARBEITEN.}
\noindent Verhältnisse und neue Probleme darbieten, denjenigen der binären Formen ganz parallel.

Zur Seite einer Form \(f\) ist stets eine zweite \(F\), ihre Adjungierte, zu betrachten, deren Koeffizienten durch diejenigen der Form \(f\) bestimmt sind, deren Determinante gleich \(\Delta^2\), und deren eigene Adjungierte im wesentlichen wieder mit \(f\) selbst identisch ist. Auch hier ist dann jede Form einer reduzierten äquivalent, in welcher ebenso wie in der dazu Adjungierten die Koeffizienten durch gewisse Ungleichheiten beschränkt sind; die Anzahl der reduzierten Formen ist endlich. Neben Zahlen sind aber durch ternäre Formen auch binäre Formen darstellbar, und zwar entspricht einer Darstellung der binären Form \(\varphi=(p,q,r)\) durch eine ternäre Form \(f\) eine Darstellung ihrer Determinante \(D\) durch deren Adjungierte \(F\) und umgekehrt; die letztern Darstellungen lassen sich also auf die erstern zurückführen. Diese aber gehören wieder zu gewissen Kongruenzwurzelpaaren \(g,\,h\ (\mathrm{mod.}\ D)\), nämlich zu einer Lösung der Kongruenzen
\[
g^2\equiv\Delta p,\ gh\equiv-\Delta q,\ h^2\equiv\Delta r\ (\mathrm{mod.}\ D)
\]
oder, wie \textsc{Gauss} kurz sagt, zur Wurzel \(\sqrt{\Delta(p,-q,r)}\ (\mathrm{mod.}\ D)\), und die Aufgabe, sie sämtlich zu finden, läuft auf die andere hinaus, die ganzzahligen Transformationen der gegebenen Form \(f\) in gewisse andere, jenen Lösungen verbundene Formen mit derselben Determinante und in sich selbst zu ermitteln.

Die Lösung der letztern Aufgabe jedoch wurde von \textsc{Gauss} nicht mehr allgemein gegeben. In den D. A. findet sie sich nur noch für die Form
\[
ax^2+a'x'^2+a''x''^2
\]
mit positiven \(a,\,a',\,a''\), und in einer kurzen, in einem am 22. November 1799 begonnenen Notizhefte (Scheda Ac, S. 22) des Nachlasses enthaltenen Notiz (W. II, p. 311) für die Form \(x^2+x'^2-x''^2\). Für Formen, durch welche nur positive Zahlen darstellbar sind, hat in der Folge ein Schüler von \textsc{Gauss}, L. A. \textsc{Seeber}, ausgehend von kristallographischen Fragestellungen, diese Aufgabe wieder aufgenommen\footnote{Siehe \textsc{Seebers} Arbeiten: \textit{Versuch einer Erklärung des innern Baues der festen Körper}, \textsc{Gilberts} Annalen der Physik \textbf{76}, 1824, S. 229 und S. 349; \textit{Untersuchung über die Eigenschaften der positiven ternären quadratischen Formen}, Freiburg im Breisgau 1831. In der ersten Arbeit wird für eine Teilung des ganzen Raumes »in kleine Parallelepipede, die sowohl einander selbst, als den parallelepipedischen Molekulen gleich sind«, das Quadrat des Abstands zweier Eckpunkte durch eine positive ternäre Form dargestellt (a. a. O. S. 353) und so eine geometrische Deutung dieser Formen gegeben. In dem Vorwort zu der Schrift von 1831 heisst es dann (S. II) in bezug auf die Aufgabe, die Kriterien für die Äquivalenz aufzustellen: »Ein Versuch, die Art, auf welche die festen Körper aus den kleinsten Teilen der Materie gebildet sind, aus den Gesetzen der Mechanik zu erklären, führte mich zur Auflösung dieser Aufgabe« und weiter (S. VIII) »Mithin ist die Theorie der positiven ternären quadratischen Formen wenigstens ein nützliches oder sogar notwendiges Hilfsmittel der Kristallographie«.}. Seine Arbeit von 1831 bedeutete aber auch darin
\newpage\SourceLeaf{34}\ArithmeticHead{27}{R}{TERNÄRE QUADRATISCHE FORMEN.\quad DER VI. ABSCHNITT DER D. A.}\HistoricalFoot{}{4*}
\noindent einen Fortschritt über \textsc{Gauss}’ Darstellung der Theorie, dass es ihm gelungen ist, die Bedingungen für reduzierte Formen so zu fassen, dass in jeder Klasse nur eine von diesen vorhanden ist, sodass die schwierige Frage nach der Äquivalenz reduzierter Formen ganz wegfällt. \textsc{Gauss} sah sich durch seine Anzeige von \textsc{Seebers} Arbeit (Gött. gel. Anz. vom 9. Juli 1831, W. II, S. 188) veranlasst, sich aufs Neue mit den ternären quadratischen Formen zu befassen, und gab zunächst den Beweis der von \textsc{Seeber} nur induktiv gemachten Bemerkung, dass für die reduzierten Formen stets \(aa'a''<2\Delta\) sei.

\textsc{Gauss} ist aber bei dieser Gelegenheit auch auf die geometrische Deutung der Theorie der quadratischen Formen, sowohl der binären wie der ternären, eingegangen, indem er ihren Zusammenhang mit den sogenannten Zahlengittern in der Ebene und im Raume, den geometrischen Sinn der eigentlichen wie der uneigentlichen Äquivalenz sowie der zuvor angeführten \textsc{Seeber}schen Ungleichheit klarlegte und im Anschluss an \textsc{Seeber} den Wert der Theorie der ternären Formen als eines Hilfsmittels für die Kristallographie hervorhob; in einem kurzen, aus dem Juli 1831 (W. II, S. 311) stammenden Fragmente (W. II, S. 305) hat \textsc{Gauss} seine bezüglichen Andeutungen durch eine Reihe raumgeometrischer Sätze weiter ausgeführt und für die letztgenannte Disziplin verwertet. Übrigens war ihm die Vorstellung der ebenen Zahlengitter, wenigstens solcher von besonderer Art, lange vor jener spätern Veröffentlichung schon vertraut. Schon die doppelte Periodizität der lemniskatischen, allgemeiner der elliptischen Funktionen, deren Untersuchung ihn bereits zu der Zeit der D. A. beschäftigte, konnte sie ihm nahe legen. Als ihm dann, mit unter dem Einflusse dieser Untersuchungen, der Gedanke kam, der Theorie der biquadratischen Reste die komplexen ganzen Zahlen \(a+bi\) zu Grunde zu legen, boten sich ihm die quadratischen, bei den kubischen Resten auch die parallelogrammatischen Zahlengitter dar, von denen er sowohl
\newpage\SourceLeaf{35}\ArithmeticHead{28}{L}{BACHMANN, ÜBER GAUSS’ ZAHLENTHEORETISCHE ARBEITEN.}
\noindent in seinen Abhandlungen W. II, S. 65 u. 93, wie in den nachgelassenen Schriftstücken W. II, S. 313—374 und VIII, S. 18 verschiedentlich Gebrauch gemacht hat; desgleichen auch bei dem geometrischen Hilfssatze W. II, S. 271, 277, auf welchen er seine Bestimmung der Klassenanzahl quadratischer Formen stützt. In der Theorie der letztern aber, wie die D. A. sie enthalten, tritt die gedachte geometrische Vorstellung nirgends zu Tage, so wenig wie die damit verwandte Zerlegbarkeit der Formen in irrationale Linearfaktoren, zu der namentlich die Lehre von der Komposition der Formen den besten Anlass bot\footnote{Es möge hier noch darauf hingewiesen werden, dass \textsc{Euler} (\textit{Vollständige Anleitung zur Algebra} II, 1770, 2. Abschnitt, 11. und 12. Kapitel, Opera omnia, Ser. I, vol. 1, S. 414 ff.) und \textsc{Lagrange} (siehe dessen \textit{Additions aux Élémens d’Algèbre d’Euler}, 1774, § IX, Oeuvres VII, S. 164) die irrationalen Linearfaktoren einer binären quadratischen Form herangezogen haben, um Fälle zu finden, unter denen eine solche Form ein Quadrat oder sonst irgend eine Potenz darstellt. \textsc{Lagrange} hat a. a. O. auch die Zusammensetzung binärer quadratischer Formen, wie überhaupt solcher Formen, die sich in Linearfaktoren zerlegen lassen, mit sich selbst bewerkstelligt, indem er von dieser Zerlegung ausgeht und die Linearfaktoren dann in abgeänderter Reihenfolge mit einander multipliziert.}; erst \textsc{Dirichlet}\begingroup\renewcommand{\thefootnote}{**)}\footnote{\textsc{Crelles} Journal f. Mathem. \textbf{19} (1839), S. 324 und \textbf{24} (1842) S. 291, Werke I, S. 411 und S. 533.}\endgroup\ hat von der letzteren ausgiebigen Gebrauch gemacht und in neuerer Zeit hat F. \textsc{Klein}\begingroup\renewcommand{\thefootnote}{***)}\footnote{Göttinger Nachrichten 1893, S. 106, ausführlicher in den Autographierten Vorlesungsheften: \textit{Ausgewählte Kapitel der Zahlentheorie} (1895/6) II 1897, S. 118.}\endgroup\ gezeigt, wie schön die Zahlengitter zur Veranschaulichung der Kompositionslehre benutzt werden können. Welche ganz neuen arithmetischen Probleme aber aus der Gittervorstellung entspringen, hat dann \textsc{Minkowski} in seiner \textit{Geometrie der Zahlen} (Leipzig, 1896) und in den bezüglichen späteren Arbeiten in weitem Umfange entwickelt.
